\subsection{Lax squares and mates}
\label{subsec:2dloc-lax-squares}

\begin{defi}
\label{def:lax-square}
A \emph{lax square} in an $(\infty,2)$-category $\A$ is a diagram
\[
\begin{tikzcd}
A \ar[r, "f"] \ar[d, "g"'] & B \ar[d, "h"] \ar[dl, Rightarrow, shorten=3mm, "\alpha"'] \\
C \ar[r, "k"'] & D
\end{tikzcd}
\]
in which the $1$-morphisms go from left to right and from top to bottom, together with a $2$-morphism $\alpha\colon hf \Rightarrow kg$ from the top-right composite to the bottom-left composite.
A \emph{commutative square} is a lax square whose $2$-morphism is invertible.
\end{defi}

The \emph{transpose} of a commutative square, obtained by reflecting it in the diagonal through $A$ and $D$, is again a commutative square, with $1$-morphisms $g$, $k$ on top and bottom, $f$, $h$ on the left and right, and $2$-morphism $\alpha^{-1}\colon kg \Rightarrow hf$.
A commutative square thus has two underlying lax squares. \textbf{Convention:} we only speak of the mates, or of the Beck--Chevalley condition, of a commutative square after its underlying lax square has been fixed explicitly.

\begin{constru}[mates]
\label{constr:mates}
Let $\alpha$ be the lax square of \cref{def:lax-square}.
\begin{enumerate}
\item Suppose that the horizontal $1$-morphisms $f$ and $k$ admit left adjoints $f^{L} \dashv f$, $k^{L} \dashv k$. Rotating the square by $90^{\circ}$ to the left, the $1$-morphisms $f$ and $k$ point upwards; replacing them by their left adjoints gives the square on the left below. Its \emph{left mate} is the lax square
\[
\begin{tikzcd}
B \ar[r, "h"] \ar[d, "f^{L}"'] & D \ar[d, "k^{L}"] \ar[dl, Rightarrow, shorten=3mm, "\lmate{\alpha}"'] \\
A \ar[r, "g"'] & C
\end{tikzcd}
\qquad
\lmate{\alpha}\colon k^{L}h \xRightarrow{k^{L}h\eta_{f}} k^{L}hff^{L} \xRightarrow{k^{L}\alpha f^{L}} k^{L}kgf^{L} \xRightarrow{\varepsilon_{k}gf^{L}} gf^{L},
\]
where $\eta_{f}\colon \mathrm{id}_{B} \Rightarrow ff^{L}$ and $\varepsilon_{k}\colon k^{L}k \Rightarrow \mathrm{id}_{C}$ are the unit and counit.
\item Suppose that the vertical $1$-morphisms $g$ and $h$ admit right adjoints $g \dashv g^{R}$, $h \dashv h^{R}$. Rotating the square by $90^{\circ}$ to the right, $g$ and $h$ point to the left; replacing them by their right adjoints, the \emph{right mate} is the lax square
\[
\begin{tikzcd}
C \ar[r, "g^{R}"] \ar[d, "k"'] & A \ar[d, "f"] \ar[dl, Rightarrow, shorten=3mm, "\rmate{\alpha}"'] \\
D \ar[r, "h^{R}"'] & B
\end{tikzcd}
\qquad
\rmate{\alpha}\colon fg^{R} \xRightarrow{\eta_{h}fg^{R}} h^{R}hfg^{R} \xRightarrow{h^{R}\alpha g^{R}} h^{R}kgg^{R} \xRightarrow{h^{R}k\varepsilon_{g}} h^{R}k,
\]
where $\eta_{h}\colon \mathrm{id}_{B} \Rightarrow h^{R}h$ and $\varepsilon_{g}\colon gg^{R} \Rightarrow \mathrm{id}_{C}$.
\end{enumerate}
Mnemonic: the left mate replaces the horizontal $1$-morphisms by their left adjoints, the right mate replaces the vertical $1$-morphisms by their right adjoints, and in both cases the result is again a lax square in the sense of \cref{def:lax-square}.
\end{constru}

\begin{defi}
\label{def:BC}
A lax square is \emph{left Beck--Chevalley} if $f$, $k$ admit left adjoints and $\lmate{\alpha}$ is invertible, and \emph{right Beck--Chevalley} if $g$, $h$ admit right adjoints and $\rmate{\alpha}$ is invertible.
\end{defi}

\begin{lem}[mate correspondence]
\label{lem:mate-correspondence}
In the situation of \cref{constr:mates}:
\begin{enumerate}
\item $\alpha \mapsto \lmate{\alpha}$ is an equivalence of spaces
$\mapp_{\Hom(A,D)}(hf,kg) \simeq \mapp_{\Hom(B,C)}(k^{L}h,gf^{L})$, and $\alpha \mapsto \rmate{\alpha}$ is an equivalence
$\mapp_{\Hom(A,D)}(hf,kg) \simeq \mapp_{\Hom(C,B)}(fg^{R},h^{R}k)$.
\item The two constructions are inverse to each other: the right mate of $\lmate{\alpha}$ (formed with the right adjoints $f$, $k$ of $f^{L}$, $k^{L}$) is $\alpha$, and the left mate of $\rmate{\alpha}$ (formed with the left adjoints $g$, $h$ of $g^{R}$, $h^{R}$) is $\alpha$.
\item Mates are compatible with horizontal and vertical pasting of lax squares, and with identity squares.
\end{enumerate}
\end{lem}

\begin{proof}
(1) Using the adjunctions of $\infty$-categories $k^{L}_{*} \dashv k_{*}$ on $\Hom(A,-)$ and $f^{*} \dashv (f^{L})^{*}$ on $\Hom(-,C)$ of \cref{subsec:2dloc-conventions}, we have equivalences
\[
\mapp(hf,kg) \simeq \mapp(k^{L}hf,g) \simeq \mapp(k^{L}h,gf^{L}),
\]
and unwinding the units and counits of these adjunctions shows that the composite is $\alpha \mapsto \lmate{\alpha}$. The right mate is treated in the same way, using $h_{*} \dashv h^{R}_{*}$ and $(g^{R})^{*} \dashv g^{*}$:
$\mapp(hf,kg) \simeq \mapp(f,h^{R}kg) \simeq \mapp(fg^{R},h^{R}k)$.
(2) and (3) are identities between composites of $2$-morphisms, which follow from the triangle identities and the interchange law in $\htwo\A$ exactly as in \cite{kellystreet1974review}.
\end{proof}

\begin{rmk}
\label{rmk:mates-memo}
In particular, a lax square whose horizontal $1$-morphisms have left adjoints and whose vertical $1$-morphisms have right adjoints has two mates, and being left and being right Beck--Chevalley are unrelated conditions. For a commutative square $\alpha$ with transpose $\alpha^{\mathrm{T}}$, the left mate of $\alpha^{\mathrm{T}}$ replaces the vertical $1$-morphisms $g$, $h$ of $\alpha$ by their left adjoints, and the right mate of $\alpha^{\mathrm{T}}$ replaces the horizontal $1$-morphisms $f$, $k$ of $\alpha$ by their right adjoints; by our convention these are mates of $\alpha^{\mathrm{T}}$, not of $\alpha$.
\end{rmk}
